\section{Introduction}
\label{sec:problem}
Verification of loose tablets and capsules is a practical challenge in medication management because removal from packaging separates a medicine from its identifying information. Appearance then provides a recognition cue, although orientation, illumination and the visible side vary, and similar colours and shapes occur across products and strengths. Within the problem of medication errors addressed by the World Health Organization's Medication Without Harm initiative~\citep{who2017medication}, visual verification concerns the association between a photographed pill and a known medicine before a medication decision is made.

A personal medication list creates an evolving recognition task in which newly prescribed medicines are enrolled from a few photographs and a query may belong to none of the enrolled identities. Recognition therefore depends on a representation that transfers across identities, together with a decision rule that distinguishes a plausible match from an unfamiliar or ambiguous pill. Since an incorrect accepted name and a request for further checking have different consequences, retrieval accuracy alone provides an incomplete account of verification performance. Existing pill retrieval benchmarks support representation learning~\citep{usuyama2020epillid,ling2020fewshot}, while the personal workflow additionally involves membership rejection and handling of close candidates.

PharmaScan formulates this task as retrieval from a personal gallery with three outcomes: an accepted identity, \textsc{unknown}, or \textsc{retake}. Medicines are registered through a shared visual embedding without network retraining, and decisions combine similarity to the nearest enrolment prototype with separation from the next candidate. Optical character recognition (OCR) using \mbox{Florence-2}~\citep{xiao2024florence2} supplies complementary imprint evidence, while catalogue retrieval provides candidates when matching within the personal gallery is insufficient. Product name and strength remain associated with the pharmacy label~\citep{moh_medication_label}, as visual similarity alone does not establish strength or suitability for administration.

The study investigates verification of previously unseen medicines from three enrolment photographs at an operating point selected on validation data, with particular attention to accepted misidentification, imprint evidence and acquisition transfer. Evaluation combines a strict identity-disjoint test from the National Library of Medicine (NLM) Computational Photography Project for Pill Identification (C3PI) collection~\citep{yaniv2016nlm}, the ePillID benchmark~\citep{usuyama2020epillid}, and MEDISEG images of Hong Kong products~\citep{mediseg2026}. Their respective protocols, detailed in Section~\ref{sec:setup}, distinguish transfer to new identities under familiar conditions from changes in cameras, product population and scene. The project plan's validation targets and corresponding measurements are reported separately from held-out results in Appendix~\ref{app:calibration}, Table~\ref{tab:targets}.

The contribution is a calibrated verification workflow and an empirical analysis of the effects of its design choices. Sequential training changes, imprint re-ranking and rejected reader adaptations are examined through accepted errors and abstention, revealing acquisition shift as a constraint even when misidentification among enrolled queries is rare. The evidence is drawn from public image datasets and simulated galleries, with patient outcomes outside the evaluation scope.
